% Options for packages loaded elsewhere
\PassOptionsToPackage{unicode}{hyperref}
\PassOptionsToPackage{hyphens}{url}
\PassOptionsToPackage{dvipsnames,svgnames,x11names}{xcolor}
%
\documentclass[
  13pt,
  a4paperpaper,
  12pt]{article}
\usepackage{amsmath,amssymb}
\usepackage{setspace}
\usepackage{iftex}
\ifPDFTeX
  \usepackage[T1]{fontenc}
  \usepackage[utf8]{inputenc}
  \usepackage{textcomp} % provide euro and other symbols
\else % if luatex or xetex
  \usepackage{unicode-math} % this also loads fontspec
  \defaultfontfeatures{Scale=MatchLowercase}
  \defaultfontfeatures[\rmfamily]{Ligatures=TeX,Scale=1}
\fi
\usepackage{lmodern}
\ifPDFTeX\else
  % xetex/luatex font selection
  \setmainfont[]{Times New Roman}
  \setmonofont[]{Courier New}
\fi
% Use upquote if available, for straight quotes in verbatim environments
\IfFileExists{upquote.sty}{\usepackage{upquote}}{}
\IfFileExists{microtype.sty}{% use microtype if available
  \usepackage[]{microtype}
  \UseMicrotypeSet[protrusion]{basicmath} % disable protrusion for tt fonts
}{}
\makeatletter
\@ifundefined{KOMAClassName}{% if non-KOMA class
  \IfFileExists{parskip.sty}{%
    \usepackage{parskip}
  }{% else
    \setlength{\parindent}{0pt}
    \setlength{\parskip}{6pt plus 2pt minus 1pt}}
}{% if KOMA class
  \KOMAoptions{parskip=half}}
\makeatother
\usepackage{xcolor}
\usepackage[margin=1.5cm,top=1.5cm,bottom=1.5cm,left=1.5cm,right=1.5cm,includeheadfoot]{geometry}
\usepackage{listings}
\newcommand{\passthrough}[1]{#1}
\lstset{defaultdialect=[5.3]Lua}
\lstset{defaultdialect=[x86masm]Assembler}
\setlength{\emergencystretch}{3em} % prevent overfull lines
\providecommand{\tightlist}{%
  \setlength{\itemsep}{0pt}\setlength{\parskip}{0pt}}
\setcounter{secnumdepth}{3}
% Minimal LaTeX preamble fragment for Pandoc-built documents
% Avoid \documentclass or loading hyperref/cleveref here (they are managed by the build pipeline)
% Provide safe unicode declarations and small custom macros only
\DeclareUnicodeCharacter{03BC}{\(\mu\)}
\DeclareUnicodeCharacter{03BB}{\(\lambda\)}
\DeclareUnicodeCharacter{03BD}{\(\nu\)}
\DeclareUnicodeCharacter{03C0}{\(\pi\)}
\DeclareUnicodeCharacter{03B5}{\(\epsilon\)}
\DeclareUnicodeCharacter{03B4}{\(\delta\)}

% Counter behaviour
\usepackage{chngcntr}
\counterwithout{figure}{section}
\counterwithout{table}{section}
\counterwithout{equation}{section}

% Caption formatting
\usepackage[font=small,labelfont=bf,textfont=it]{caption}
\captionsetup{justification=raggedright,singlelinecheck=false}
\captionsetup[figure]{position=bottom,skip=10pt}
\captionsetup[table]{position=top,skip=10pt}
\AtBeginEnvironment{figure}{\raggedright}

% Titlesec formatting
\usepackage{titlesec}
\titlespacing{\section}{0pt}{12pt}{6pt}
\titlespacing{\subsection}{0pt}{10pt}{4pt}
\titlespacing{\subsubsection}{0pt}{8pt}{3pt}
\titleformat{\section}{\large\bfseries}{\thesection}{1em}{}
\titleformat{\subsection}{\normalsize\bfseries}{\thesubsection}{1em}{}
\titleformat{\subsubsection}{\normalsize\itshape}{\thesubsubsection}{1em}{}

% Graphics support for figures
\usepackage{graphicx}

% Small visual tweaks
\usepackage{xcolor}
\definecolor{codebg}{RGB}{245,245,245}
\definecolor{codeborder}{RGB}{200,200,200}
\definecolor{codefg}{RGB}{50,50,50}

% Setup listings safely (don't override global styles)
\usepackage{listings}
\lstset{backgroundcolor=\color{codebg},basicstyle=\normalsize\ttfamily\color{codefg},frame=single,framerule=0.5pt,framesep=5pt}

% Keep title metadata to be inserted by build system
\ifLuaTeX
  \usepackage{selnolig}  % disable illegal ligatures
\fi
\IfFileExists{bookmark.sty}{\usepackage{bookmark}}{\usepackage{hyperref}}
\IfFileExists{xurl.sty}{\usepackage{xurl}}{} % add URL line breaks if available
\urlstyle{same}
\hypersetup{
  colorlinks=true,
  linkcolor={blue},
  filecolor={blue},
  citecolor={blue},
  urlcolor={blue},
  pdfcreator={LaTeX via pandoc}}

\author{}
\date{}

\begin{document}

{
\hypersetup{linkcolor=black}
\setcounter{tocdepth}{3}
\tableofcontents
}
\setstretch{1.0}
\section{Ant Stack Implementation
Appendix}\label{sec:ant_stack_appendix}

\subsection{Introduction}\label{introduction}

This appendix demonstrates the practical implementation of the
vibrational theory of olfaction within the Ant Stack cognitive
architecture framework. The Ant Stack provides a structured three-layer
approach (AntBody, AntBrain, AntMind) for modeling insect intelligence,
and we map our computational modules to this architecture using thin
adapter patterns.

The implementation emphasizes: - \textbf{Scientific accuracy}: Direct
integration with validated \passthrough{\lstinline!src/!} computational
utilities - \textbf{Behavioral realism}: Incorporating empirical
constraints from the literature - \textbf{Reproducibility}: Fixed random
seeds and deterministic algorithms - \textbf{Extensibility}: Modular
design enabling species-specific customization

This mapping serves as both a validation of our theoretical framework
and a practical tool for behavioral simulation and hypothesis testing.

\subsubsection{AntBody Layer: Physical Simulation and
Sensing}\label{antbody-layer-physical-simulation-and-sensing}

\paragraph{Sensilla Morphology
Integration}\label{sensilla-morphology-integration}

\begin{lstlisting}[language=Python]
# AntBody sensilla configuration (adapter pattern)
class AntBodySensilla:
    def __init__(self, species_preset: str):
        # Load species-specific sensilla parameters via cohereAnts presets
        self.lengths = load_sensilla_lengths(species_preset)
        self.diameters = load_sensilla_diameters(species_preset)
        # Delegate resonance calculation to tested src utilities
        from src.sensilla import calculate_wavelength_matching
        self.optimal_wavelengths = calculate_wavelength_matching(self.lengths, self.diameters)

    def export_io(self) -> dict:
        return {
            'lengths_um': self.lengths,
            'diameters_um': self.diameters,
            'optimal_wavelengths_um': self.optimal_wavelengths,
        }
\end{lstlisting}

\textbf{I/O Contract}: - \textbf{Observations}: Sensilla dimensions
(μm), resonance frequencies (THz), quality factors - \textbf{Actions}:
Antenna positioning, sensilla orientation - \textbf{Physics}: 1 kHz
update rate, contact dynamics for substrate interaction

\paragraph{Spectroscopy and Atmospheric
Transmission}\label{spectroscopy-and-atmospheric-transmission}

Integration of cohereAnts atmospheric transmission models:

\begin{lstlisting}[language=Python]
class AntBodySpectroscopy:
    def __init__(self, environment_preset: str):
        self.transmission_curves = load_atmospheric_data(environment_preset)
        self.spectral_resolution = 0.01  # μm
        
    def get_transmission(self, wavelength: float, distance: float) -> float:
        # Delegate to cohereAnts atmospheric transmission model in src/core
        return calculate_atmospheric_transmission(wavelength, distance)
\end{lstlisting}

\textbf{Configuration Parameters}: - Atmospheric windows: 2-5 μm, 8-14
μm, 17-25 μm - Transmission coefficients: 0.7-0.9 for optimal windows -
Distance-dependent attenuation models

\subsubsection{AntBrain Layer: Neural
Architecture}\label{antbrain-layer-neural-architecture}

\paragraph{Olfactory Processing
Pipeline}\label{olfactory-processing-pipeline}

Mapping cohereAnts vibrational theory to AntBrain's AL→MB→CX
architecture:

\begin{lstlisting}[language=Python]
class AntBrainOlfaction:
    def __init__(self, neuron_count: int = 100000):
        # Antennal Lobe (AL) - odor coding
        self.al_neurons = self._initialize_al_circuit()
        # Mushroom Body (MB) - associative learning
        self.mb_neurons = self._initialize_mb_circuit()
        # Central Complex (CX) - spatial integration
        self.cx_neurons = self._initialize_cx_circuit()
    
    def _initialize_al_circuit(self):
        # Delegate vibrational detection to src components in production
        # Each glomerulus responds to specific molecular vibrations
        return VibrationalGlomeruliCircuit()
    
    def _initialize_mb_circuit(self):
        # Kenyon cells for odor-memory associations
        return KenyonCellCircuit()
    
    def _initialize_cx_circuit(self):
        # Ring attractor for heading representation
        return RingAttractorCircuit()
\end{lstlisting}

\textbf{Neural Implementation Details}: - \textbf{AL Layer}: 50
glomeruli, each tuned to specific vibrational frequencies - \textbf{MB
Layer}: 2500 Kenyon cells with sparse coding (5\% activity) - \textbf{CX
Layer}: 16-heading ring attractor with 100 neurons per heading

\paragraph{Vibrational Detection
Circuit}\label{vibrational-detection-circuit}

Implementation of cohereAnts electromagnetic theory:

\begin{lstlisting}[language=Python]
class VibrationalGlomeruliCircuit:
    def __init__(self):
        self.frequency_tuning = np.linspace(2, 25, 50)  # μm to THz
        self.quality_factors = np.ones(50) * 100
        
    def process_spectral_input(self, spectral_data: np.ndarray) -> np.ndarray:
        # Implement cohereAnts resonance detection
        responses = np.zeros(50)
        for i, freq in enumerate(self.frequency_tuning):
            responses[i] = self._calculate_vibrational_response(spectral_data, freq)
        return responses
    
    def _calculate_vibrational_response(self, spectrum: np.ndarray, 
                                     resonant_freq: float) -> float:
        # Placeholder: call src electromagnetic coupling utilities in production
        coupling_strength = self._calculate_coupling(spectrum, resonant_freq)
        return coupling_strength * self.quality_factors[i]
\end{lstlisting}

\subsubsection{AntMind Layer: Cognitive
Modeling}\label{antmind-layer-cognitive-modeling}

\paragraph{Active Inference for Olfactory
Search}\label{active-inference-for-olfactory-search}

Integration of cohereAnts behavioral models with active inference:

\begin{lstlisting}[language=Python]
class AntMindOlfaction:
    def __init__(self):
        self.generative_model = self._build_olfactory_model()
        self.policy_horizon = 2.0  # seconds
        
    def _build_olfactory_model(self):
        # Implement cohereAnts behavioral predictions
        return OlfactoryGenerativeModel()
    
    def select_policy(self, current_state: Dict) -> np.ndarray:
        # Active inference policy selection
        expected_free_energy = self._calculate_efe()
        return self._minimize_free_energy(expected_free_energy)
    
    def _calculate_efe(self) -> Dict[str, float]:
        # Decompose into epistemic and pragmatic value
        return {
            'epistemic': self._calculate_epistemic_value(),
            'pragmatic': self._calculate_pragmatic_value()
        }
\end{lstlisting}

\paragraph{Stigmergy for Trail
Following}\label{stigmergy-for-trail-following}

Implementation of cohereAnts pheromone dynamics:

\begin{lstlisting}[language=Python]
class AntMindStigmergy:
    def __init__(self):
        self.pheromone_field = np.zeros((100, 100))
        self.decay_rate = 0.01
        self.diffusion_coefficient = 0.1
        
    def update_pheromone_field(self, deposits: List[Tuple[int, int, float]]):
        # Implement cohereAnts pheromone diffusion model
        for x, y, amount in deposits:
            self.pheromone_field[x, y] += amount
        
        # Apply diffusion and decay
        self.pheromone_field = self._diffuse_and_decay()
    
    def _diffuse_and_decay(self) -> np.ndarray:
        # Fick's law implementation from cohereAnts
        laplacian = self._calculate_laplacian(self.pheromone_field)
        diffusion = self.diffusion_coefficient * laplacian
        decay = -self.decay_rate * self.pheromone_field
        return self.pheromone_field + diffusion + decay
\end{lstlisting}

\subsection{Species-Specific
Implementations}\label{species-specific-implementations}

\subsubsection{Formica Species
Configuration}\label{formica-species-configuration}

\begin{lstlisting}[language=Python]
# Formica species preset for Ant Stack
FORMICA_PRESET = {
    'body': {
        'sensilla_lengths': [15.2, 18.7, 22.1, 19.8, 16.5],  # μm
        'sensilla_diameters': [2.1, 2.8, 3.2, 2.9, 2.3],     # μm
        'optimal_wavelengths': [60.8, 74.8, 88.4, 79.2, 66.0], # μm
        'antenna_length': 2.5,  # mm
        'leg_count': 6,
        'body_mass': 0.015  # g
    },
    'brain': {
        'al_glomeruli_count': 50,
        'mb_kenyon_cells': 2500,
        'cx_heading_resolution': 16,
        'spiking_threshold': 0.1,
        'learning_rate': 0.01
    },
    'mind': {
        'policy_horizon': 2.0,  # seconds
        'pheromone_decay': 0.01,
        'diffusion_coefficient': 0.1,
        'exploration_rate': 0.2
    }
}
\end{lstlisting}

\subsubsection{Camponotus Species
Configuration}\label{camponotus-species-configuration}

\begin{lstlisting}[language=Python]
# Camponotus species preset for Ant Stack
CAMPONOTUS_PRESET = {
    'body': {
        'sensilla_lengths': [22.5, 28.1, 31.7, 26.8, 24.3],  # μm
        'sensilla_diameters': [3.2, 4.1, 4.8, 4.2, 3.6],     # μm
        'optimal_wavelengths': [90.0, 112.4, 126.8, 107.2, 97.2], # μm
        'antenna_length': 3.8,  # mm
        'leg_count': 6,
        'body_mass': 0.045  # g
    },
    'brain': {
        'al_glomeruli_count': 60,
        'mb_kenyon_cells': 3000,
        'cx_heading_resolution': 20,
        'spiking_threshold': 0.08,
        'learning_rate': 0.015
    },
    'mind': {
        'policy_horizon': 2.5,  # seconds
        'pheromone_decay': 0.008,
        'diffusion_coefficient': 0.12,
        'exploration_rate': 0.15
    }
}
\end{lstlisting}

\subsection{Evaluation and
Benchmarking}\label{evaluation-and-benchmarking}

\subsubsection{Navigation Performance
Metrics}\label{navigation-performance-metrics}

\begin{lstlisting}[language=Python]
class AntStackEvaluator:
    def __init__(self, test_scenarios: List[str]):
        self.scenarios = test_scenarios
        self.metrics = {}
    
    def evaluate_navigation(self, ant_stack: AntStack) -> Dict[str, float]:
        results = {}
        for scenario in self.scenarios:
            if scenario == 'trail_following':
                results[scenario] = self._evaluate_trail_following(ant_stack)
            elif scenario == 'food_search':
                results[scenario] = self._evaluate_food_search(ant_stack)
            elif scenario == 'nest_return':
                results[scenario] = self._evaluate_nest_return(ant_stack)
        return results
    
    def _evaluate_trail_following(self, ant_stack: AntStack) -> float:
        # Implement cohereAnts trail following metrics (calls src/behavioral metrics)
        trail_deviation = self._calculate_trail_deviation()
        pheromone_detection = self._calculate_pheromone_detection()
        return self._combine_metrics([trail_deviation, pheromone_detection])
    
    def _evaluate_food_search(self, ant_stack: AntStack) -> float:
        # Implement cohereAnts search efficiency metrics (calls src/behavioral metrics)
        search_time = self._measure_search_time()
        energy_efficiency = self._calculate_energy_efficiency()
        return self._combine_metrics([search_time, energy_efficiency])
\end{lstlisting}

\subsubsection{Robustness Testing}\label{robustness-testing}

\begin{lstlisting}[language=Python]
class RobustnessTester:
    def __init__(self):
        self.noise_levels = [0.01, 0.05, 0.1, 0.2]
        self.adversary_types = ['sensor_noise', 'pheromone_contamination', 'path_obstruction']
    
    def test_noise_robustness(self, ant_stack: AntStack) -> Dict[str, float]:
        results = {}
        for noise_level in self.noise_levels:
            performance = self._run_noisy_scenario(ant_stack, noise_level)
            results[f'noise_{noise_level}'] = performance
        return results
    
    def test_adversary_robustness(self, ant_stack: AntStack) -> Dict[str, float]:
        results = {}
        for adversary in self.adversary_types:
            performance = self._run_adversarial_scenario(ant_stack, adversary)
            results[f'adversary_{adversary}'] = performance
        return results
\end{lstlisting}

\subsection{Implementation Workflow}\label{implementation-workflow}

\subsubsection{Development Pipeline}\label{development-pipeline}

\begin{enumerate}
\def\labelenumi{\arabic{enumi}.}
\tightlist
\item
  \textbf{Module Mapping}: Identify cohereAnts functions for Ant Stack
  integration
\item
  \textbf{I/O Contract Definition}: Establish standardized interfaces
  between layers
\item
  \textbf{Species Preset Creation}: Develop parameterized configurations
\item
  \textbf{Testing Framework}: Implement evaluation metrics and
  benchmarks
\item
  \textbf{Documentation}: Create implementation guides and examples
\end{enumerate}

\subsubsection{Code Organization}\label{code-organization}

\begin{lstlisting}
ant_stack_cohereants/
├── antbody/
│   ├── sensilla_physics.py      # cohereAnts vibrational theory
│   ├── spectroscopy_sensors.py  # atmospheric transmission models
│   └── morphology_models.py     # species-specific parameters
├── antbrain/
│   ├── olfactory_circuits.py    # AL→MB→CX implementation
│   ├── vibrational_detection.py # electromagnetic coupling
│   └── learning_mechanisms.py   # STDP and plasticity
├── antmind/
│   ├── olfactory_inference.py   # active inference models
│   ├── stigmergy_models.py      # pheromone dynamics
│   └── behavioral_policies.py   # search and navigation
├── presets/
│   ├── formica_config.py        # Formica species preset
│   ├── camponotus_config.py     # Camponotus species preset
│   └── custom_species.py        # Template for new species
└── evaluation/
    ├── navigation_tests.py      # Trail following, search
    ├── robustness_tests.py      # Noise, adversary testing
    └── performance_metrics.py   # Standardized benchmarks
\end{lstlisting}

\subsection{Integration Benefits}\label{integration-benefits}

\subsubsection{Reproducibility}\label{reproducibility}

\begin{itemize}
\tightlist
\item
  \textbf{Standardized I/O}: All experiments use consistent interfaces
\item
  \textbf{Version Pinning}: Dependencies and parameters are explicitly
  tracked
\item
  \textbf{Seed Management}: Reproducible random number generation
\item
  \textbf{Artifact Tracking}: Complete experiment provenance
\end{itemize}

\subsubsection{Extensibility}\label{extensibility}

\begin{itemize}
\tightlist
\item
  \textbf{Species Presets}: Easy addition of new ant species
\item
  \textbf{Module Swapping}: Interchangeable components across layers
\item
  \textbf{Parameter Tuning}: Systematic exploration of parameter space
\item
  \textbf{Benchmark Addition}: New evaluation scenarios
\end{itemize}

\subsubsection{Validation}\label{validation}

\begin{itemize}
\tightlist
\item
  \textbf{Biological Plausibility}: Grounded in empirical data
\item
  \textbf{Performance Metrics}: Quantified success criteria
\item
  \textbf{Robustness Testing}: Resilience to real-world challenges
\item
  \textbf{Cross-Species Transfer}: Generalization across taxa
\end{itemize}

\subsection{Future Directions}\label{future-directions}

\subsubsection{Advanced Learning
Mechanisms}\label{advanced-learning-mechanisms}

\begin{itemize}
\tightlist
\item
  \textbf{Meta-Learning}: Adaptation across different environments
\item
  \textbf{Collective Intelligence}: Emergent behaviors in colonies
\item
  \textbf{Transfer Learning}: Knowledge transfer between species
\end{itemize}

\subsubsection{Hardware Integration}\label{hardware-integration}

\begin{itemize}
\tightlist
\item
  \textbf{Robotic Platforms}: Physical ant-inspired robots
\item
  \textbf{Sensor Networks}: Distributed environmental monitoring
\item
  \textbf{Edge Computing}: Efficient on-device processing
\end{itemize}

\subsubsection{Biological Validation}\label{biological-validation}

\begin{itemize}
\tightlist
\item
  \textbf{Field Studies}: Comparison with natural ant behavior
\item
  \textbf{Neural Recording}: Validation against biological data
\item
  \textbf{Evolutionary Analysis}: Phylogenetic patterns in behavior
\end{itemize}

\subsection{Conclusion}\label{conclusion}

The integration of cohereAnts research into the Ant Stack framework
provides a robust, reproducible platform for studying ant intelligence.
By mapping our vibrational theory of olfaction, spectroscopic analysis,
and behavioral modeling to the standardized three-layer architecture, we
create a comprehensive system that bridges theoretical insights with
computational implementation.

This implementation enables systematic exploration of ant behavior
across species, environments, and experimental conditions while
maintaining the biological plausibility that underpins our research. The
modular design facilitates both hypothesis testing in myrmecology and
applications in swarm robotics, cognitive security, and AI alignment.

\textbf{Key Contributions}: 1. \textbf{Systematic Integration}:
Methodical mapping of cohereAnts to Ant Stack layers 2. \textbf{Species
Parameterization}: Reproducible configurations for multiple ant taxa 3.
\textbf{Evaluation Framework}: Standardized metrics and robustness
testing 4. \textbf{Implementation Workflow}: Clear development pipeline
and code organization 5. \textbf{Future Roadmap}: Extensibility and
validation pathways

The resulting framework serves as a bridge between theoretical
entomology and computational neuroscience, enabling reproducible
research that advances our understanding of both natural ant
intelligence and artificial intelligence systems.

\end{document}
